\documentclass[11pt]{article}
\usepackage{labsheet_common}

\title{}
\date{}

\begin{document}
\thispagestyle{fancy}

\labtitle{AI Training -- Lab 2}{Detecting Rail Defects with AI}{A hands-on introduction to image segmentation and model tuning}
\nocodebadge

You will follow copy-paste steps in a terminal and edit one plain-text
settings file. That's it -- the AI code itself is already written for you.

\section*{What you'll learn}
\begin{itemize}
  \item What ``Non-Destructive Testing'' (NDT) means and why AI can help with it
  \item What it means for AI to ``segment'' an image (draw where something is, pixel by pixel) instead of just labeling the whole photo
  \item How a single setting -- the \textbf{decision threshold} -- controls the trade-off between catching every defect and avoiding false alarms
  \item How to read an accuracy score (IoU) and compare it with classmates
\end{itemize}

\section*{The big picture}
Railway maintenance crews inspect rail surfaces for defects (cracks,
spalling, corrosion) that could cause problems if left unnoticed --
traditionally by a human walking the track and looking. NDT means checking
for defects without damaging the rail (e.g. just photographing it), which is
exactly what our dataset does.

\begin{notebox}
Unlike Lab 1, \textbf{nobody needs to label anything here} -- the dataset
already comes with 113 real rail photos, each with a hand-drawn ``ground
truth'' map showing exactly where the defect is (drawn by the researchers
who published this dataset). Your job is to train an AI that draws its own
version of that map on a new photo, and to make that AI as accurate as
possible.
\end{notebox}

\section*{What's already done for you}
\begin{itemize}
  \item The dataset (113 photos + their official defect maps) has already been downloaded into \path|datasets/raw/rail_ndt/rsdds/|.
  \item All the AI code -- feature extraction, model training, scoring, plotting -- is already written in \texttt{scripts/}. You will only \emph{run} it, never edit it.
  \item The only file you'll edit is \texttt{configs/rail\_ndt\_config.yaml}, a plain-text settings file (not code).
\end{itemize}

\section*{How it all fits together}
\begin{center}
\resizebox{\textwidth}{!}{%
\begin{tikzpicture}[node distance=6mm]
  \node[autobox] (data) {RSDDs photos + ground truth (already downloaded)};
  \node[actionbox, right=of data] (config) {Edit\\\texttt{rail\_ndt\_\allowbreak config.yaml}};
  \node[autobox, right=of config] (train) {Train model\\\texttt{rail\_\allowbreak train.py}};
  \node[autobox, right=of train] (evaluate) {Evaluate model\\\texttt{rail\_\allowbreak evaluate.py}};
  \node[scorebox, right=of evaluate] (score) {Your score:\\Mean IoU};

  \draw[flowarrow] (data) -- (config);
  \draw[flowarrow] (config) -- (train);
  \draw[flowarrow] (train) -- (evaluate);
  \draw[flowarrow] (evaluate) -- (score);

  \draw[loopB] (score.south) to[out=-55, in=-125, looseness=0.8]
    node[looplabel, midway, below=1mm] {just change threshold -- no retrain (Part 3)} (evaluate.south);
  \draw[loopA] (score.north) to[out=55, in=125, looseness=0.9]
    node[looplabel, midway, above=1mm] {change model settings -- needs retrain (Part 4)} (config.north);
\end{tikzpicture}%
}
\end{center}
\diagramlegend

Two different loops back into the same pipeline: the dashed blue one under
the diagram is fast (just re-run evaluate); the solid orange one over the
top requires starting again from training.

\section*{Part 0: One-time setup}
Open a terminal (on Mac: the \textbf{Terminal} app; on Windows:
\textbf{Anaconda Prompt} or \textbf{PowerShell}) and run these commands one
at a time:

\begin{lstlisting}
cd AI-training
python3 -m venv .venv
source .venv/bin/activate      # Windows: .venv\Scripts\activate
pip install -r requirements.txt
\end{lstlisting}

If you already did this for Lab 1 in the same folder, you can skip straight
to Part 1 (just remember to run the \texttt{source .venv/bin/activate} line
again in any new terminal window).

\section*{Part 1: Train your model}
Open \texttt{configs/rail\_ndt\_config.yaml} in any plain-text editor
(Notepad, TextEdit, VS Code -- anything that isn't Word). You'll see
settings like:

\begin{lstlisting}
model_type: random_forest
n_estimators: 100
max_depth: 8
random_seed: 42
threshold: 0.5
\end{lstlisting}

For now, leave everything as-is and run:

\begin{lstlisting}
python3 scripts/rail_train.py
\end{lstlisting}

This looks at every pixel of the training photos and learns what a ``defect
pixel'' tends to look like (its brightness, texture, and where in the photo
it is), compared to normal rail surface. It takes under a minute.

\section*{Part 2: Evaluate your model}
\begin{lstlisting}
python3 scripts/rail_evaluate.py
\end{lstlisting}

This tests your model on photos it never trained on, and prints something
like:

\begin{lstlisting}
Mean IoU:  0.25 (higher is better, 1.0 = perfect)
Mean Dice: 0.38 (higher is better, 1.0 = perfect)
\end{lstlisting}

\textbf{IoU (Intersection over Union)} measures how much your model's
guessed defect area overlaps with the real defect area. 1.0 = perfect
overlap, 0 = no overlap at all. This is your score -- \textbf{higher is
better}.

It also saves a picture, \texttt{results/figures/rail\_eval.png}, showing
several test photos side by side: the original photo, the real defect map,
and your model's guessed defect map. Look at it -- this is the most useful
way to understand what your model is actually doing.

\section*{Part 3: Tune the threshold (the main lever -- no retraining needed)}
Your model doesn't just say ``defect'' or ``not defect'' for each pixel --
it gives a confidence score from 0 to 1. The \texttt{threshold} setting
decides how confident it must be before marking a pixel as a defect:

\begin{center}
\begin{tabular}{@{}lp{9.5cm}@{}}
\toprule
\textbf{threshold} & \textbf{Effect} \\
\midrule
Low (e.g. 0.2)  & Flags more pixels -- catches more real defects, but also more false alarms \\
High (e.g. 0.8) & Flags fewer pixels -- fewer false alarms, but risks missing real defects \\
\bottomrule
\end{tabular}
\end{center}

Open \texttt{configs/rail\_ndt\_config.yaml}, change only the
\texttt{threshold} value, and re-run \textbf{just} the evaluate step (no
need to retrain):

\begin{lstlisting}
python3 scripts/rail_evaluate.py
\end{lstlisting}

\begin{tipbox}
Try several values (e.g. 0.3, 0.4, 0.5, 0.6, 0.7, 0.8) and note the Mean IoU
each time. This is the fastest way to improve your score, and needs no
retraining -- each check takes only a few seconds.
\end{tipbox}

\section*{Part 4: (Optional/stretch) Retrain with different settings}
If you want to go further, edit \texttt{model\_type}, \texttt{n\_estimators},
\texttt{max\_depth}, or \texttt{random\_seed}, then re-run \textbf{both}
steps in order (these settings only take effect after retraining):

\begin{lstlisting}
python3 scripts/rail_train.py
python3 scripts/rail_evaluate.py
\end{lstlisting}

Then repeat Part 3 to re-tune the threshold for your new model -- the best
threshold can be different for different settings.

\section*{Compete with classmates}
Keep a scoreboard:

\begin{center}
\small
\begin{tabular}{@{}*{5}{>{\raggedright\arraybackslash}p{2.9cm}}@{}}
\toprule
\textbf{Your name} & \textbf{model\_type} & \textbf{key settings} & \textbf{threshold} & \textbf{Mean IoU (higher=better)} \\
\midrule
 & & & & \\[10pt]
 & & & & \\[10pt]
 & & & & \\[10pt]
\bottomrule
\end{tabular}
\end{center}

\section*{Glossary}
\begin{itemize}
  \item \textbf{NDT (Non-Destructive Testing)} -- inspecting something for defects without damaging it, e.g. from a photo instead of cutting the rail open.
  \item \textbf{Segmentation} -- instead of labeling a whole photo (``defective'' or not), the model marks \emph{which pixels} belong to the defect.
  \item \textbf{Ground truth} -- the correct answer, here a hand-drawn map of exactly where each defect is, made by the dataset's original researchers.
  \item \textbf{Threshold} -- the confidence cutoff (0 to 1) above which a pixel is called ``defect.'' The only setting you can change without retraining.
  \item \textbf{IoU (Intersection over Union)} -- overlap between your model's guessed defect area and the real one, divided by their combined area. 1.0 = identical, 0 = no overlap.
  \item \textbf{Dice score} -- a very similar overlap measure to IoU; usually a bit higher for the same prediction. Both are reported so you can compare.
  \item \textbf{False positive} -- a normal pixel wrongly marked as a defect (a false alarm).
  \item \textbf{False negative} -- a real defect pixel the model missed.
  \item \textbf{Class imbalance} -- in every photo, the vast majority of pixels are normal rail and only a small area is defect. The model is specifically built to account for this, otherwise it could get a deceptively high score just by never predicting ``defect'' at all.
  \item \textbf{Random Forest} -- a model made of many simple decision-making ``trees'' that vote on each pixel; \texttt{n\_estimators} = number of trees, \texttt{max\_depth} = how detailed each tree's questions can get.
\end{itemize}

\begin{warnbox}
\begin{itemize}
  \item \textbf{``no RSDDs images found''} -- run \texttt{python3 scripts/fetch\_rsdds.py} first to download the dataset (only needs to be done once).
  \item \textbf{``no trained model found''} -- run \texttt{python3 scripts/rail\_train.py} before \texttt{rail\_evaluate.py}.
  \item \textbf{Changed a setting but the score didn't change} -- only \texttt{threshold} takes effect immediately; every other setting requires re-running \texttt{rail\_train.py} first.
  \item \textbf{Score looks much worse than a classmate's with identical settings} -- double check \texttt{random\_seed} is the same; it controls which photos land in the training vs. test group, so a different seed is training/testing on a different split of photos entirely.
\end{itemize}
\end{warnbox}

\section*{Something to think about (discuss with the class)}
\begin{itemize}
  \item Why is a high IoU score harder to achieve here than a high accuracy score would be, given how small the defect areas are compared to the whole photo?
  \item In a real inspection system, would you rather have a lower threshold (more false alarms) or a higher one (more missed defects)? Does your answer change if a missed crack could cause a derailment?
  \item This dataset has no ``clean rail'' photos without any defect at all -- every image contains one. How might that affect what the model has learned, and what would you want to add to fix it?
\end{itemize}

\end{document}
